By definition, identification over \(\mathfrak F_x(C,\varepsilon)\) means that \(\theta_{y\leftarrow x}(\mathcal M)\) has the same value for every completion in that fiber. The sharp-effect-frontier theorem identifies its image with the nonempty set \(\mathcal R_{xy}(C)\). Therefore identification is equivalent to \(|\mathcal R_{xy}(C)|=1\).

For every \(j\in J_x(C)\), \(C_{xj}\ne0\). Hence the row of \(C_{\{x,y\},J_x(C)}\) indexed by \(x\) is nonzero. This two-row matrix has rank one exactly when every two of its columns have zero determinant, that is,
\[
 C_{xj}C_{yk}-C_{xk}C_{yj}=0
 \quad\text{for all }j,k\in J_x(C). \tag{12}
\]
Division by the nonzero \(C_{xj}C_{xk}\) shows that (12) is equivalent to
\[
 \frac{C_{yj}}{C_{xj}}=\frac{C_{yk}}{C_{xk}}
 \quad\text{for all }j,k\in J_x(C),
\]
which is exactly the singleton condition.

If the rank condition fails, choose \(j,k\in J_x(C)\) for which the cross-product in (12) is nonzero. Their two ratios are distinct. The witness algorithm returns same-law completions \(\mathcal M_j\) and \(\mathcal M_k\), and the intervention-ratio lemma gives the two distinct effects. After the shared \(O(n^3)\) basis extension and inverse, each of the two complementary matchings and row normalizations costs \(O(n^3)\), so the additional arithmetic cost remains \(O(n^3)\).